\documentclass[../../main.tex]{subfiles}% 注意这里的文件路径不能用 ./main.tex ,否则用latexmk编译子文件会报错
\graphicspath{{\subfix{./image/}}} % 指定图片目录,后续可以直接使用图片文件名
% 注意这里的文件路径不能用 ../../image/ ,否则用latexmk编译子文件会报错

% 例如:
% \begin{figure}[H]
% \centering
% \includegraphics[scale=0.3]{图.png}
% \caption{}
% \label{figure:图}
% \end{figure}
% 注意:上述\label{}一定要放在\caption{}之后,否则引用图片序号会只会显示??.

\begin{document}

\section{完备化}

\begin{definition}
\begin{enumerate}
\item 设 $(\mathscr{X},\rho),(\mathscr{X}_1,\rho_1)$ 是两个度量空间, 如果存在映射 $\varphi:\mathscr{X}\to\mathscr{X}_1$ 满足
\begin{enumerate}[(1)]
\item $\varphi$ 是满射,
\item $\rho(x,y)=\rho_1(\varphi x,\varphi y)\quad (\forall x,y\in\mathscr{X})$,
\end{enumerate}
则称 $(\mathscr{X},\rho)$ 和 $(\mathscr{X}_1,\rho_1)$ 是\textbf{等距同构的}, 并称 $\varphi$ 为\textbf{等距同构映射}, 有时简称\textbf{等距同构}.

\item 如果度量空间 $(\mathscr{X}_1,\rho_1)$ 与另一个度量空间 $(\mathscr{X}_2,\rho_2)$之间存在映射$\varphi:\mathscr{X}_1\to \mathscr{X}_2$,满足
\begin{align*}
\rho_1(x,y)=\rho_2(\varphi x,\varphi y)\quad (\forall x,y\in\mathscr{X}_1).
\end{align*}
也即$(\mathscr{X}_1,\rho_1)$ 与$(\mathscr{X}_2,\rho_2)$ 的子空间 $(\varphi(\mathscr{X}_1),\rho_2)$ 是等距同构的.
此时我们称$(\mathscr{X}_1,\rho_1)$ \textbf{可以嵌入} $(\mathscr{X}_2,\rho_2)$. 在等距同构意义下, 我们认为 $(\mathscr{X}_1,\rho_1)$ 就是 $(\mathscr{X}_2,\rho_2)$ 的一个\textbf{子空间}, 并简单地记作
$$
(\mathscr{X}_1,\rho_1)\subseteq(\mathscr{X}_2,\rho_2).
$$
\end{enumerate}
\end{definition}
\begin{remark}
由(2)容易导出 $\varphi$ 还是单射.
\end{remark}
\begin{remark}
凡是等距同构的度量空间, 它们的一切与距离相联系的性质都是一样的. 因此今后我们将不再区分它们. 
\end{remark}

\begin{definition}[稠密子集]
设 $(\mathscr{X},\rho)$ 是度量空间. 集合 $E\subseteq\mathscr{X}$ 叫作在 $\mathscr{X}$ 中的\textbf{稠密子集}, 如果 $\forall x\in\mathscr{X},\forall\varepsilon>0,\exists z\in E$ 使得 $\rho(x,z)<\varepsilon$. 换句话说: $\forall x\in\mathscr{X},\exists\{x_n\}\subseteq E$ 使得 $x_n\to x\ (n\to\infty)$.
\end{definition}

\begin{proposition}[稠密子集的传递性]\label{proposition:稠密子集的传递性}
设 $(X,\rho)$ 是一个度量空间，$A \subseteq B \subseteq X$。如果 $A$ 在 $B$ 中稠密，且 $B$ 在 $X$ 中稠密，那么 $A$ 在 $X$ 中稠密。
\end{proposition}
\begin{proof}
任取 $x\in X$ 和 $\varepsilon>0$。由 $B$ 在 $X$ 中稠密，存在 $b\in B$ 使得
\[
\rho(x,b) < \frac{\varepsilon}{2}.
\]
由于 $b\in B$，且 $A$ 在 $B$ 中稠密，存在 $a\in A$ 使得
\[
\rho(b,a) < \frac{\varepsilon}{2}.
\]
根据三角不等式，
\[
\rho(x,a) \le \rho(x,b) + \rho(b,a) < \frac{\varepsilon}{2} + \frac{\varepsilon}{2} = \varepsilon.
\]
因此，对任意 $x\in X$ 和 $\varepsilon>0$，均存在 $a\in A$ 满足 $\rho(x,a)<\varepsilon$，即 $A$ 在 $X$ 中稠密。
\end{proof}

\begin{definition}[完备化空间]
包含给定度量空间 $(\mathscr{X},\rho)$ 的最小的完备度量空间称为 $\mathscr{X}$ 的\textbf{完备化空间}, 其中最小的含义是: 任何一个以 $(\mathscr{X},\rho)$ 为子空间的完备度量空间都以此空间为子空间.
\end{definition}

\begin{proposition}\label{proposition:稠密的度量子空间得到完备化空间}
如果 $(\mathscr{X}_1,\rho_1)$ 是一个以 $(\mathscr{X},\rho)$ 为度量子空间的完备度量空间, $\rho_1|_{\mathscr{X}\times\mathscr{X}}=\rho$, 并且 $\mathscr{X}$ 在 $\mathscr{X}_1$ 中稠密, 则 $\mathscr{X}_1$ 是 $\mathscr{X}$ 的完备化空间.
\end{proposition}
\begin{proof}
由$\mathscr{X}$ 在 $\mathscr{X}_1$ 中稠密知,对$\forall\xi\in\mathscr{X}_1,\exists x_n\in\mathscr{X}$ 使得 $\rho_1(x_n,\xi)\to0\ (n\to\infty)$. 如果存在完备度量空间$(\mathscr{X}_2,\rho_2)$ 以 $(\mathscr{X},\rho)$ 为度量子空间, 因为
\begin{align*}
\rho_2(x_n,x_m)=\rho_1(x_n,x_m)\to0\quad (n,m\to\infty).
\end{align*}
因此$\{x_n\}$在$\mathscr{X}_2$中是基本列,所以由$(\mathscr{X}_2,\rho_2)$的完备性知$\exists\hat{\xi}\in\mathscr{X}_2$ 使得 $\rho_2(x_n,\hat{\xi})\to0$. 做映射 $T:\mathscr{X}_1\to\mathscr{X}_2,T\xi=\hat{\xi}$. 我们只需证 $T$ 是等距的. 因为 $\forall\eta\in\mathscr{X}_1$, 又 $\exists y_n\in\mathscr{X}$ 使得 $\rho_1(y_n,\eta)\to0(n\to \infty)$, 所以
\begin{align*}
\rho_1(\xi,\eta) \leqslant \lim\limits_{n\rightarrow \infty} \left[ \rho_1(\xi,x_n) + \rho_1(x_n,y_n) + \rho_1(y_n,\eta) \right] = \lim\limits_{n\rightarrow \infty} \rho_1(x_n,y_n),\\
\lim\limits_{n\rightarrow \infty} \rho_1(x_n,y_n) \leqslant \lim\limits_{n\rightarrow \infty} \left[ \rho_1(x_n,\xi) + \rho_1(\xi,\eta) + \rho_1(\eta,y_n) \right] = \rho_1(\xi,\eta),
\end{align*}
从而$\rho_1(\xi,\eta) = \lim\limits_{n\rightarrow \infty} \rho_1(x_n,y_n).$ 同理可得$\rho_2(\hat{\xi},\hat{\eta}) = \lim\limits_{n\rightarrow \infty} \rho_2(x_n,y_n).$ 又因为$\{x_n\},\{y_n\} \subseteq \mathscr{X} \subseteq \mathscr{X}_1\cap \mathscr{X}_2,$ 所以
\begin{align*}
\rho_1(x_n,y_n) = \rho(x_n,y_n) = \rho_2(x_n,y_n),\quad \forall n\in \mathbb{N}.
\end{align*}
于是
\begin{align*}
\rho_1(\xi,\eta) = \lim\limits_{n\rightarrow \infty} \rho_1(x_n,y_n) = \lim\limits_{n\rightarrow \infty} \rho_2(x_n,y_n) = \rho_2(\hat{\xi},\hat{\eta}).
\end{align*}
这表明 $(\mathscr{X}_1,\rho_1)$ 可以嵌入 $(\mathscr{X}_2,\rho_2)$. 在等距同构意义下,$(\mathscr{X}_1,\rho_1)$ 是$(\mathscr{X}_2,\rho_2)$的子空间.

\end{proof}

\begin{theorem}\label{theorem:每一个度量空间都有一个完备化空间}
每一个度量空间都有一个完备化空间.
\end{theorem}
\begin{proof}
设 $(\mathscr{X},\rho)$ 是一个度量空间, 分三步证明它有一个完备化空间.
\begin{enumerate}[(1)]
\item 将 $\mathscr{X}$ 中的基本列分类, 凡是满足
\begin{align*}
\lim\limits_{n\to\infty}\rho(x_n,y_n)=0
\end{align*}
的两个基本列 $\{x_n\},\{y_n\}$ 称为等价的. 彼此等价的基本列归于同一类且只归一类, 称为等价类. 我们把一个等价类看成是一个元素, 并用 $\mathscr{X}_1$ 表示一切这种元素(等价类)组成的集合. 在 $\mathscr{X}_1$ 上定义距离: $\forall\xi,\eta\in\mathscr{X}_1$, 任取 $\{x_n\}\in\xi,\{y_n\}\in\eta$, 令
\begin{align}
\rho_1(\xi,\eta)=\lim\limits_{n\to\infty}\rho(x_n,y_n).\label{eq:1.2.2}
\end{align}
容易验证 \eqref{eq:1.2.2} 式右端的极限的确存在并且极限值与 $\{x_n\},\{y_n\}$ 的选取无关. 为了验证 \eqref{eq:1.2.2} 式定义的 $\rho_1$ 的确是个距离, 注意到定义中的非负性、唯一性和交换性是显然的, 设$\{x_n\},\{y_n\},\{z_n\}$ 是分别属于等价类 $\xi,\eta,\zeta$ 的基本列,则对
\begin{align*}
\rho(x_n,y_n)\leqslant\rho(x_n,z_n)+\rho(z_n,y_n)
\end{align*}
取极限得到三角不等式
\begin{align*}
\rho(\xi,\eta)\leqslant\rho(\xi,\zeta)+\rho(\zeta,\eta).
\end{align*}
这样, 我们就证明了 $(\mathscr{X}_1,\rho_1)$ 是一度量空间.

\item $\forall x\in\mathscr{X}$, 我们用 $\xi_x\in\mathscr{X}_1$ 表示包含序列 $(x,x,\cdots,x,\cdots)$ 的等价类, 这样的 $\xi_x$ 全体记为 $\mathscr{X}'$. 显然 $\mathscr{X}'\subseteq\mathscr{X}_1$ 且映射 $T:x\mapsto\xi_x$ 作为 $(\mathscr{X},\rho)\to(\mathscr{X}',\rho_1)$ 的映射显是等距同构映射. 因此, $(\mathscr{X},\rho)$ 和 $(\mathscr{X}',\rho_1)$ 等距同构, 即有 $(\mathscr{X},\rho)\subseteq(\mathscr{X}_1,\rho_1)$. 进一步容易验证 $\mathscr{X}$ 在 $\mathscr{X}_1$ 中稠密.

\item 证明 $(\mathscr{X}_1,\rho_1)$ 是完备的. 设 $\{\xi^{(n)}\}$ 是 $\mathscr{X}_1$ 中的基本列. 要证 $\exists\xi\in\mathscr{X}_1$ 使得
\begin{align*}
\rho_1(\xi^{(n)},\xi)\to0\quad (n\to\infty).
\end{align*}
\begin{enumerate}[(i)]
\item 先证特殊情形, 假定 $\{\xi^{(n)}\}\subseteq\mathscr{X}'$. 令 $x_n=T^{-1}\xi^{(n)}$, 则 $\{x_n\}$ 是 $\mathscr{X}$ 中的基本列. 设 $\{x_n\}\in\xi\in\mathscr{X}_1$, 便有 $\xi^{(n)}\to\xi$.

\item 再证一般情形, 由于 $\mathscr{X}'$ 在 $\mathscr{X}_1$ 中稠密, $\forall\xi^{(n)}\in\mathscr{X}_1$, $\exists\overline{\xi}^{(n)}\in\mathscr{X}'$ 使得 $\rho_1(\overline{\xi}^{(n)},\xi^{(n)})<\frac{1}{n}$. 由(i)可设 $\overline{\xi}^{(n)}\to\xi\in\mathscr{X}_1$, 即可推出 $\xi^{(n)}\to\xi$.
\end{enumerate}
\end{enumerate}
最后综合(1)(2)(3)并用\refpro{proposition:稠密的度量子空间得到完备化空间}即得结论.

\end{proof}

\begin{proposition}[常见度量空间的完备化空间]\label{proposition:常见度量空间的完备化空间}
\begin{enumerate}[(1)]
\item $P[a,b]$ ($[a,b]$ 上的多项式全体)在 $C[a,b]$ 中稠密.定义距离
$$\rho(x,y)\triangleq \max_{a\leqslant t\leqslant b}|x(t)-y(t)|,$$
则$(P[a,b],\rho)$的完备化空间是 $(C[a,b],\rho)$.

\item 定义距离
\begin{align*}
\rho _1(x,y)\triangleq \int_a^b{\left| x\left( t \right) -y\left( t \right) \right|\mathrm{d}t},
\end{align*}
则$(C[a,b],\rho_1)$ 是不完备的度量空间, 其完备化空间是 $(L^1[a,b],\rho_1)$.
\end{enumerate}
\end{proposition}
\begin{proof}
\begin{enumerate}[(1)]
\item 根据\hyperref[Basis of Analytics-theorem:多项式逼近有限阶光滑函数]{Weierstrass定理}可立得$P[a,b]$在 $C[a,b]$ 中稠密. 由\refthe{theorem:常见的度量空间}知$(C[a,b],\rho)$是完备的度量空间.显然$(P[a,b],\rho)$是其度量子空间,再由\refpro{proposition:稠密的度量子空间得到完备化空间}可知$(C[a,b],\rho)$就是$(P[a,b],\rho)$的完备化空间.

\item 显然$(C[a,b],\rho_1)$是度量空间。取
\begin{align*}
f(t)=
\begin{cases}
0,\quad 0\leqslant t\leqslant \frac{1}{2},\\
1,\quad \frac{1}{2}\leqslant t\leqslant 1.
\end{cases},\quad f_n(t)=
\begin{cases}
0,\quad 0\leqslant t\leqslant \frac{1}{2}-\frac{1}{n},\\
n\left(t-\frac{1}{2}+\frac{1}{n}\right),\quad \frac{1}{2}-\frac{1}{n}\leqslant t\leqslant \frac{1}{2},\\
1,\quad \frac{1}{2}\leqslant t\leqslant 1,
\end{cases}
\,\, \forall n\in\mathbb{N},
\end{align*}
则$\{f_n\}\subseteq C[a,b]$，且
\begin{align*}
\rho_1(f_n,f)\to 0\ (n\to\infty)\Longleftrightarrow f_n\to f\ (n\to\infty).
\end{align*}
但$f(t)$在$t=\frac{1}{2}$处不连续，故$f\notin C[0,1]$。因此$(C[a,b],\rho_1)$是不完备的度量空间。
由\refthe{Real Analysis-theorem:可积函数可被有紧支集的函数逼近}知$(C[a,b],\rho_1)$在$(L^1[a,b],\rho_1)$中稠密，故由\refpro{proposition:稠密的度量子空间得到完备化空间}知$(L^1[a,b],\rho_1)$是$(C[a,b],\rho_1)$的完备化空间。
\end{enumerate}

\end{proof}

\begin{proposition}\label{proposition:空间S是完备度量空间}
令$S$为一切实(或复)数列
\begin{align*}
x=(\xi_1,\xi_2,\cdots,\xi_n,\cdots)
\end{align*}
组成的集合, 在$S$中定义距离为
\begin{align*}
\rho(x,y)=\sum\limits_{k=1}^{\infty}\frac{1}{2^k}\cdot\frac{|\xi_k-\eta_k|}{1+|\xi_k-\eta_k|},
\end{align*}
其中$x=(\xi_1,\xi_2,\cdots,\xi_k,\cdots), y=(\eta_1,\eta_2,\cdots,\eta_k,\cdots)$。则 $(S,\rho)$为一个完备的度量空间.称$(S,\rho)$是\textbf{空间}$\bm{S}$.
\end{proposition}
\begin{proof}
显然$\rho$有正定性。对$\forall x=(\xi_1,\xi_2,\cdots,\xi_k,\cdots),y=(\eta_1,\eta_2,\cdots,\eta_k,\cdots),z=(\zeta_1,\zeta_2,\cdots,\zeta_k,\cdots)\in S$，则
\begin{align*}
\rho(x,y)&=\sum\limits_{k=1}^{\infty}\frac{1}{2^k}\cdot\frac{|\xi_k-\eta_k|}{1+|\xi_k-\eta_k|}\leqslant\sum\limits_{k=1}^{\infty}\frac{1}{2^k}\cdot\frac{|\xi_k-\zeta_k|+|\zeta_k-\eta_k|}{1+|\xi_k-\zeta_k|+|\zeta_k-\eta_k|}\\
&=\sum\limits_{k=1}^{\infty}\frac{1}{2^k}\cdot\frac{|\xi_k-\zeta_k|}{1+|\xi_k-\zeta_k|+|\zeta_k-\eta_k|}+\sum\limits_{k=1}^{\infty}\frac{1}{2^k}\cdot\frac{|\zeta_k-\eta_k|}{1+|\xi_k-\zeta_k|+|\zeta_k-\eta_k|}\\
&\leqslant\sum\limits_{k=1}^{\infty}\frac{1}{2^k}\cdot\frac{|\xi_k-\zeta_k|}{1+|\xi_k-\zeta_k|}+\sum\limits_{k=1}^{\infty}\frac{1}{2^k}\cdot\frac{|\zeta_k-\eta_k|}{1+|\zeta_k-\eta_k|}\\
&=\rho(x,z)+\rho(z,y).
\end{align*}
故$(S,\rho)$是一个度量空间。
再设$\{x^{(n)}\}_{n=1}^{\infty}$是$S$中的Cauchy列，其中$x^{(n)}=(\xi_1^{(n)},\xi_2^{(n)},\cdots,\xi_k^{(n)},\cdots),\forall n\in\mathbb{N}$。则对$\forall n,m\in\mathbb{N}$且$m\geqslant n$，有
\begin{align*}
&\frac{1}{2^k}\cdot\frac{|\xi_k^{(m)}-\xi_k^{(n)}|}{1+|\xi_k^{(m)}-\xi_k^{(n)}|}\leqslant\sum\limits_{k=1}^{\infty}\frac{1}{2^k}\cdot\frac{|\xi_k^{(m)}-\xi_k^{(n)}|}{1+|\xi_k^{(m)}-\xi_k^{(n)}|}=\rho(x^{(m)},x^{(n)})=0,\forall k\in\mathbb{N}\\
\implies&\lim\limits_{n\rightarrow\infty}\left(\frac{1}{2^k}\cdot\frac{|\xi_k^{(m)}-\xi_k^{(n)}|}{1+|\xi_k^{(m)}-\xi_k^{(n)}|}\right)\leqslant\lim\limits_{n\rightarrow\infty}\rho(x^{(m)},x^{(n)})=0,\forall k\in\mathbb{N}\\
\implies&\lim\limits_{n\rightarrow\infty}|\xi_k^{(m)}-\xi_k^{(n)}|=0,\forall k\in\mathbb{N}.
\end{align*}
因此对$\forall k\in\mathbb{N}$，$\{\xi_k^{(n)}\}_{n=1}^{\infty}\subseteq(\mathbb{R},d)$都是Cauchy列，其中
\begin{align*}
d\triangleq|x-y|,\quad\forall x,y\in\mathbb{R}.
\end{align*}
又注意到$(\mathbb{R},d)$是完备度量空间，故存在$\xi_k'\in\mathbb{R}$，使得$\lim\limits_{n\rightarrow\infty}d(\xi_k^{(n)},\xi_k')=\lim\limits_{n\rightarrow\infty}|\xi_k^{(n)}-\xi_k'|=0$。令$x'=(\xi_1',\xi_2',\cdots,\xi_k',\cdots)$，则由
\begin{align*}
\left|\frac{1}{2^k}\cdot\frac{|\xi_k^{(n)}-\xi_k'|}{1+|\xi_k^{(n)}-\xi_k'|}\right|\leqslant\frac{1}{2^k},\forall k\in\mathbb{N};\quad\sum\limits_{k=1}^{\infty}\frac{1}{2^k}<+\infty;\quad\lim\limits_{n\rightarrow\infty}\frac{|\xi_k^{(n)}-\xi_k'|}{1+|\xi_k^{(n)}-\xi_k'|}=0,
\end{align*}
以及\hyperref[Basis of Analytics-theorem:级数的控制收敛定理]{级数的控制收敛定理}得
\begin{align*}
\lim\limits_{n\rightarrow\infty}\rho(x^{(n)},x')=\lim\limits_{n\rightarrow\infty}\sum\limits_{k=1}^{\infty}\frac{1}{2^k}\cdot\frac{|\xi_k^{(n)}-\xi_k'|}{1+|\xi_k^{(n)}-\xi_k'|}=\sum\limits_{k=1}^{\infty}\frac{1}{2^k}\cdot\lim\limits_{n\rightarrow\infty}\frac{|\xi_k^{(n)}-\xi_k'|}{1+|\xi_k^{(n)}-\xi_k'|}=0.
\end{align*}
因此$(S,\rho)$是一个完备度量空间。

\end{proof}

\begin{proposition}\label{proposition:度量空间中基本列收敛当且仅当其子列收敛}
在一个度量空间$(X,\rho)$上, 求证: 基本列是收敛列当且仅当这个基本列中存在一串收敛子列.
并且这个收敛子列的极限就是原基本列的极限.
\end{proposition}
\begin{proof}
必要性显然成立，下证充分性。设$\{x_n\}$是$X$中的基本列，且存在一串收敛子列$\{x_{n_k}\}$，不妨设$\lim\limits_{k\rightarrow \infty}\rho(x_{n_k},x_0)=0$，其中$x_0\in X$，则$n_k\geqslant k,\forall k\in\mathbb{N}$。于是由$\{x_n\}$是基本列知
\begin{align*}
\lim\limits_{k\rightarrow \infty}\rho(x_k,x_{n_k})=0.
\end{align*}
又因为对$\forall k\in\mathbb{N}$，有
\begin{align*}
\rho(x_k,x_0)\leqslant \rho(x_k,x_{n_k})+\rho(x_{n_k},x_0).
\end{align*}
所以对上式令$k\rightarrow \infty$，得
\begin{align*}
\lim\limits_{k\rightarrow \infty}\rho(x_k,x_0)\leqslant \lim\limits_{k\rightarrow \infty}\rho(x_k,x_{n_k})+\lim\limits_{k\rightarrow \infty}\rho(x_{n_k},x_0)=0\Longrightarrow \lim\limits_{k\rightarrow \infty}\rho(x_k,x_0)=0.
\end{align*}
故基本列$\{x_n\}$是收敛列。

\end{proof}

\begin{example}
设$F$是只有有限项不为$0$的实数列全体, 在$F$上引进距离
\begin{align*}
\rho(x,y)=\sup\limits_{k\geqslant1}|\xi_k-\eta_k|,
\end{align*}
其中$x=\{\xi_k\}\in F, y=\{\eta_k\}\in F$。求证: $(F,\rho)$是不完备的度量空间, 并指出它的完备化空间.
\end{example}
\begin{proof}
显然$(F,\rho)$是度量空间。取
\begin{align*}
x^{(n)}=\left\{x_k^{(n)}\right\}=\left\{\underbrace{1,\frac{1}{2},\frac{1}{3},\cdots,\frac{1}{n}}_{\text{前}n\text{项}},0,0\cdots\right\},\forall n\in\mathbb{N}.
\end{align*}
则$\{x^{(n)}\}\subseteq F$。对$\forall n,m\in\mathbb{N}$且$m\geqslant n$，有
\begin{align*}
\left\{x_k^{(m)}-x_k^{(n)}\right\}=\left\{\underbrace{\overbrace{0,\cdots,0}^{\text{前}n\text{项}},\frac{1}{n+1},\cdots,\frac{1}{m}}_{\text{前}m\text{项}},0,0,\cdots\right\}\Longrightarrow\lim\limits_{n\rightarrow\infty}\rho(x^{(m)},x^{(n)})=\lim\limits_{n\rightarrow\infty}\frac{1}{n+1}=0.
\end{align*}
故$\{x^{(n)}\}$是基本列。再令$x=\{x_k\}=\left\{1,\frac{1}{2},\cdots,\frac{1}{n},\cdots\right\}\notin F$（每一项都非零），则
\begin{align*}
\left\{x_k^{(m)}-x_k\right\}=\left\{\underbrace{0,\cdots,0}_{\text{前}n\text{项}},\frac{1}{n+1},\frac{1}{n+2},\cdots\right\}\Longrightarrow\lim\limits_{n\rightarrow\infty}\rho(x^{(n)},x)=\lim\limits_{n\rightarrow\infty}\frac{1}{n+1}=0.
\end{align*}
故$\{x^{(n)}\}$收敛到$x$，但$x\notin F$。因此$(F,\rho)$是不完备度量空间。

记$l_0^\infty$是极限为$0$的实数列全体，先证$(l_0^\infty,\rho)$是完备度量空间。

设$\{x^{(n)}\}_{n=1}^{\infty}$是$l_0^\infty$中的基本列，其中$x^{(n)}=\{x_k^{(n)}\}_{k=1}^{\infty}$。则对$\forall \varepsilon>0$，存在$N\in\mathbb{N}$，使得对$\forall n,m\in\mathbb{N}$且$m\geqslant n>N$，有
\begin{align*}
\left|x_k^{(n)}-x_k^{(m)}\right|\leqslant\sup\limits_{k\geqslant 1}\left|x_k^{(n)}-x_k^{(m)}\right|=\rho(x^{(n)},x^{(m)})<\varepsilon,\quad \forall k\in\mathbb{N}.
\end{align*}
于是对$\forall k\in\mathbb{N}$，$\{x_k^{(n)}\}_{n=1}^{\infty}$是$(\mathbb{R},d)$中的基本列，其中$d\triangleq|x-y|,\forall x,y\in\mathbb{R}$。显然$(\mathbb{R},d)$是完备度量空间，
故对$\forall k\in\mathbb{N}$，$\{x_k^{(n)}\}_{n=1}^{\infty}$都是收敛列。即对$\forall k\in\mathbb{N}$，存在$x_k^0\in\mathbb{R}$，使得$\lim\limits_{m\rightarrow\infty}\left|x_k^{(m)}-x_k^0\right|=0$。

令$x_0=(x_1^0,x_2^0,\cdots,x_k^0,\cdots)$，则对$\forall n,m\in\mathbb{N}$且$m\geqslant n>N$，有
\begin{align*}
\left|x_k^{(n)}-x_k^0\right|\leqslant\left|x_k^{(n)}-x_k^{(m)}\right|+\left|x_k^{(m)}-x_k^0\right|\leqslant\sup\limits_{k\geqslant 1}\left|x_k^{(n)}-x_k^{(m)}\right|+\left|x_k^{(m)}-x_k^0\right|<\varepsilon+\left|x_k^{(m)}-x_k^0\right|,\quad \forall k\in\mathbb{N}.
\end{align*}
令$m\rightarrow\infty$得，对$\forall n>N$有
\begin{align}\label{eq:FHJIOWJIOafesf3u0rfwf}
\left|x_k^{(n)}-x_k^0\right|<\varepsilon,\quad \forall k\in\mathbb{N}.
\end{align}
于是对$\forall n>N$有
\begin{align*}
\rho(x^{(n)},x_0)=\sup\limits_{k\geqslant 1}\left|x_k^{(n)}-x_k^0\right|\leqslant\varepsilon.
\end{align*}
故$\lim\limits_{n\rightarrow\infty}\rho(x^{(n)},x_0)=0$。由$\{x^{(n)}\}\in l_0^\infty$知$\lim\limits_{k\rightarrow\infty}x_k^{(N+1)}=0$，再结合\eqref{eq:FHJIOWJIOafesf3u0rfwf}式知，存在$K\in\mathbb{N}$，使得$\forall k>K$，有
\begin{align*}
\left|x_k^0\right|\leqslant\left|x_k^{(N+1)}-x_k^0\right|+\left|x_k^{(N+1)}\right|<2\varepsilon.
\end{align*}
即$\lim\limits_{k\rightarrow\infty}x_k^0=0$，因此$x_0\in l_0^\infty$。故$(l_0^\infty,\rho)$是完备度量空间。

再证$F$在$l_0^\infty$中稠密。对$\forall x=\{\xi_1,\xi_2,\cdots,\xi_k,\cdots\}\in l_0^\infty$，则$\lim\limits_{k\rightarrow\infty}\xi_k=0$。取
\begin{align*}
x^{(n)}=\left\{\underbrace{\xi_1,\xi_2,\cdots,\xi_n}_{\text{前}n\text{项}},0,0\cdots\right\},\forall n\in\mathbb{N}.
\end{align*}
则$\{x^{(n)}\}_{n=1}^{\infty}\subseteq F$，并且
\begin{align*}
\rho(x^{(n)},x)=\sup\limits_{k\geqslant 1}\left\{\xi_{n+1},\xi_{n+2},\cdots\right\}=\sup\limits_{k\geqslant n+1}\xi_k,\quad \forall n\in\mathbb{N}.
\end{align*}
令$n\rightarrow\infty$得
\begin{align*}
\lim\limits_{n\rightarrow\infty}\rho(x^{(n)},x)=\lim\limits_{n\rightarrow\infty}\sup\limits_{k\geqslant n+1}\xi_k=\varlimsup\limits_{k\rightarrow\infty}\xi_k=0.
\end{align*}
故$F$在$l_0^\infty$中稠密。
综上,由\refpro{proposition:稠密的度量子空间得到完备化空间}知$(l_0^\infty,\rho)$是$(F,\rho)$的完备化空间。

\end{proof}

\begin{example}
求证: $[0,1]$上的多项式全体$P[0,1]$按距离
\begin{align*}
\rho(p,q)=\int_{0}^{1}|p(x)-q(x)|\mathrm{d}x \quad (p,q\text{ 是多项式})
\end{align*}
是不完备的, 并指出它的完备化空间.
\end{example}
\begin{proof}
显然$(P[0,1],\rho)$是度量空间。
取$f(x)=\left|x-\frac{1}{2}\right|$，则由\rrefthe{Basis of Analytics-theorem:可积被连续函数逼近}{Basis of Analytics-theorem:可积被连续函数逼近-2}知，存在多项式序列$\{p_n\}_{n=1}^{\infty}$，使得
\begin{align*}
\lim\limits_{n\rightarrow \infty}\int_0^1\left|f(x)-p_n(x)\right|\mathrm{d}x=0.
\end{align*}
从而对$\forall n,m\in\mathbb{N}$且$m\geqslant n$，有
\begin{align*}
\rho (p_m,p_n)=\int_0^1{\left| p_m\left( x \right) -p_n\left( x \right) \right|\mathrm{d}x}\leqslant \int_0^1{\left| p_m\left( x \right) -f\left( x \right) \right|\mathrm{d}x}+\int_0^1{\left| f\left( x \right) -p_n\left( x \right) \right|\mathrm{d}x}.
\end{align*}
令$n\rightarrow \infty$得$\lim\limits_{n\rightarrow \infty}\rho(p_m,p_n)=0$，故$\{p_n\}_{n=1}^{\infty}$是$(P[0,1],\rho)$中的基本列。但其收敛函数$f(x)\notin P[0,1]$，因此$\{p_n\}_{n=1}^{\infty}$不是$(P[0,1],\rho)$中的收敛列，故$(P[0,1],\rho)$是不完备的。

由\refthe{Real Analysis-theorem:Lp完备的距离空间}知$(L^1[0,1],\rho)$是完备度量空间。由\refthe{Real Analysis-theorem:可积函数可被有紧支集的函数逼近}知$C[0,1]$在$L^1[0,1]$中稠密。
又由\rrefthe{Basis of Analytics-theorem:可积被连续函数逼近}{Basis of Analytics-theorem:可积被连续函数逼近-2}
知$P[0,1]$在$C[0,1]$中稠密，故由\hyperref[proposition:稠密子集的传递性]{稠密子集的传递性}知$P[0,1]$在$L^1[0,1]$中稠密。
因此由\refpro{proposition:稠密的度量子空间得到完备化空间}知$(L^1[0,1],\rho)$是$(P[0,1],\rho)$的完备化空间。
\end{proof}

\begin{example}
在完备的度量空间$(\mathscr{X},\rho)$中给定点列$\{x_n\}$, 如果$\forall \varepsilon>0$, 存在基本列$\{y_n\}$, 使得
\begin{align*}
\rho(x_n,y_n)<\varepsilon \quad (n\in\mathbb{N}),
\end{align*}
求证: $\{x_n\}$收敛.
\end{example}
\begin{proof}
由条件知，对$\forall \varepsilon>0$，存在基本列$\{y_n\}$，使得
\begin{align*}
\rho(x_n,y_n)<\varepsilon,\quad \forall n\in\mathbb{N}.
\end{align*}
从而存在$N\in\mathbb{N}$，使得对$\forall n,m\in\mathbb{N}$且$m\geqslant n>N$，有
\begin{align*}
\rho(y_n,y_m)<\varepsilon.
\end{align*}
于是对$\forall n,m\in\mathbb{N}$且$m\geqslant n>N$，有
\begin{align*}
\rho(x_n,x_m)\leqslant\rho(x_n,y_n)+\rho(y_n,y_m)+\rho(y_m,x_m)<3\varepsilon.
\end{align*}
故$\{x_n\}$是基本列。又因为$(\mathscr{X},\rho)$是完备的，所以$\{x_n\}$是收敛列。

\end{proof}
































\end{document}